\honest{The Proper Use of Humility}{Eliezer Yudkowsky}{2006}

Everyone agrees that good science needs humility. The question is what sort. I sort some cases. The creationist who tells you to be humble about evolution is practicing ``a very selective underconfidence.'' The engineer who builds fail-safes while ``damn sure'' the machine will not fail is humble in the good way, and so is the student who double-checks a test. The student who never checks, because checking cannot bring certainty, is less wise.

Then there is the student who says ``one so lowly as myself can hope for no better lot.'' ``This is social modesty, not humility.'' Social modesty is about status in the tribe, an ``ancestrally relevant concept.'' Scientific humility is ``a more recent and rarefied invention,'' which you would practice alone in a spacesuit. I give no evidence for either history. \nb{The dictionary's senses fit the lowly student better than they fit double-checking. Merriam-Webster defines humility as ``freedom from pride or arrogance,'' and humble as ``not proud or haughty'' and ``ranking low in a hierarchy or scale.'' Double-checking and fail-safes are nearer to what it calls prudence. A proud person can take precautions; the post's engineer is ``damn sure.''}

When the student cites the problem of induction, I answer with a motive: the student wants ``to hand in the test quickly, and go home and play video games.''

Physicists, with ``this arrogant idea'' that their models should always work, follow up small flaws, and now and then a flaw overturns a theory. ``Therefore it is written: `If you do not seek perfection you will halt before taking your first steps.'\,'' Then I turn to satire. Think of the social audacity of trying to be right all the time. If science said evolution was true most of the time, or that the Earth might be flat on some days, scientists would seem less confrontational. They have won a little status with medicine and cellphones and now think themselves chiefs of the tribe, so they should ``compromise a little.''

A principle you understand only vaguely can be bent to justify almost anything. People with vague models end up believing what they wanted to, but ``the purpose of our ethics is to move us, not be moved by us.'' ``The vast majority of appeals that I witness'' to humility are excuses to shrug: the lottery player who says you cannot know he will lose, the doubter of evolution who says you cannot prove it, the person who calls a hard problem probably too hard. The cause is motivated skepticism: scrutinizing more heavily what we do not want to believe. Humility misunderstood becomes a fully general excuse, since you can never be sure. \nb{The cited study, by Taber and Lodge, defines it as the post says.}

Humility is also easy to profess, like the religious ``belief in belief'' that I credit to Daniel Dennett. It is easy to answer every argument with ``Well, of course I could be wrong'' and then change nothing. As Galbraith said, faced with changing one's mind or proving there is no need to, ``almost everyone gets busy on the proof.'' I have often seen people accept new information and then explain why they will do exactly what they planned before, for a new reason. The point of thinking is to shape plans, and humility misunderstood is a wonderful black hole for news. The test is whether the actions a humble line of thought recommends make you ``stronger, or weaker.'' A bridge needs extra cables, not a shrug.

``Therefore it is written: `To be humble is to take specific actions in anticipation of your own errors. To confess your fallibility and then do nothing about it is not humble; it is boasting of your modesty.'\,'' \nb{Both sentences introduced by ``Therefore it is written'' come from Yudkowsky's own ``Twelve Virtues of Rationality.'' The post does not say so.}
